% =============================================================================
% Appendix D: Glossary
% =============================================================================
\chapter{Glossary}
\label{app:glossary}

\begin{description}[style=nextline, leftmargin=2cm]

\item[Attention]
Mechanism that allows each token to ``look at'' other tokens in the sequence. Computes weighted combinations based on relevance. See Chapter~\ref{ch:self-attention}.

\item[Attention Mask]
Binary tensor indicating which positions should be attended to (1) or ignored (0). Used to prevent attending to padding tokens.

\item[Autoregressive]
Generation strategy where each token is predicted based only on previous tokens. The model generates one token at a time, feeding each output back as input.

\item[Batch Size]
Number of training examples processed together in one forward/backward pass. Larger batches provide more stable gradients but require more memory.

\item[BOS (Beginning of Sequence)]
Special token marking the start of a sequence. In \TryLM: \texttt{<|bos|>}.

\item[BPE (Byte Pair Encoding)]
Tokenization algorithm that iteratively merges the most frequent character pairs. Creates a vocabulary of subword units.

\item[Causal Masking]
Attention pattern where each position can only attend to earlier positions. Prevents the model from ``seeing the future'' during training.

\item[Checkpoint]
Saved state of a model during training, including weights, optimizer state, and configuration. Enables resuming training and inference.

\item[Context Length]
Maximum sequence length the model can process. \TryLM uses 512 tokens.

\item[Cross-Entropy Loss]
Loss function measuring the difference between predicted probability distribution and true labels. Standard loss for classification tasks including language modeling.

\item[d\_model]
Dimension of token embeddings and hidden states throughout the model. \TryLM uses 384.

\item[Decoder-Only]
Transformer architecture using only the decoder portion with causal masking. Used for text generation (GPT-style). Contrasts with encoder-only (BERT) or encoder-decoder (T5).

\item[Dropout]
Regularization technique that randomly sets activations to zero during training. Prevents overfitting.

\item[Embedding]
Dense vector representation of a discrete token. Maps vocabulary indices to continuous vectors.

\item[EOS (End of Sequence)]
Special token marking the end of a sequence. Signals generation should stop. In \TryLM: \texttt{<|eos|>}.

\item[Epoch]
One complete pass through the entire training dataset.

\item[Feed-Forward Network (FFN)]
Two-layer neural network applied to each token independently. In transformers, typically expands dimension then projects back.

\item[Fine-tuning]
Training a pre-trained model on a specific task or dataset. Adapts general knowledge to specialized domains.

\item[Flash Attention]
Optimized attention implementation achieving O(n) memory instead of O(n²). Uses tiling to avoid materializing the full attention matrix.

\item[FLOPs]
Floating-point operations. Measure of computational cost. One forward+backward pass costs approximately 6N FLOPs per token, where N is parameter count.

\item[Gradient Accumulation]
Technique to simulate larger batch sizes by accumulating gradients over multiple forward passes before updating weights.

\item[Gradient Clipping]
Limiting gradient magnitude to prevent exploding gradients. Typically clip to max norm of 1.0.

\item[Head (Attention Head)]
One of multiple parallel attention computations in multi-head attention. Each head can learn different attention patterns.

\item[Hidden Size]
See d\_model.

\item[Inference]
Using a trained model to make predictions on new data. Contrast with training.

\item[KV Cache]
Stored key and value tensors from previous positions during generation. Avoids recomputing attention for already-generated tokens.

\item[Language Model]
Model that predicts probability distributions over tokens. Can be used for text generation, completion, and understanding.

\item[LayerNorm]
Normalization technique that normalizes across features for each example. Stabilizes training.

\item[Learning Rate]
Step size for optimizer updates. Controls how much weights change each step. Typical values: 1e-4 to 1e-3.

\item[Logits]
Raw (unnormalized) output scores from the model before softmax. Shape: (batch, sequence, vocab\_size).

\item[Loss]
Scalar value measuring how wrong the model's predictions are. Training minimizes this value.

\item[Mixed Precision]
Training with both float16 and float32 precision. Reduces memory and increases speed while maintaining accuracy.

\item[MoE (Mixture of Experts)]
Architecture using multiple ``expert'' FFN layers with a router selecting which experts process each token.

\item[Multi-Head Attention]
Attention mechanism with multiple parallel attention heads, each operating on a portion of the embedding dimension.

\item[Next-Token Prediction]
The core task of language models: given a sequence, predict the probability distribution over possible next tokens.

\item[Optimizer]
Algorithm that updates model weights based on gradients. AdamW is standard for transformers.

\item[Overfitting]
When a model performs well on training data but poorly on new data. Sign: training loss decreases while validation loss increases.

\item[PAD (Padding)]
Special token used to fill sequences to equal length for batching. Usually masked in attention. In \TryLM: \texttt{<|pad|>}.

\item[Parameters]
Learnable weights in the model. \TryLM has approximately 10.8M parameters.

\item[Perplexity]
Evaluation metric for language models. Exponential of cross-entropy loss: $\text{PPL} = \exp(\text{loss})$. Lower is better.

\item[Position Encoding]
Information added to embeddings indicating token position in the sequence. Can be learned, sinusoidal, or rotary (RoPE).

\item[Pre-norm]
Transformer variant where normalization is applied before (rather than after) attention and FFN layers. More stable training.

\item[Pre-training]
Initial training phase on large amounts of unlabeled text. Learns general language patterns.

\item[Query, Key, Value (Q, K, V)]
Three projections of the input used in attention. Query asks ``what am I looking for?'', Key says ``what do I contain?'', Value provides ``what information to pass.''

\item[Residual Connection]
Adding the input of a layer to its output: $x + f(x)$. Helps gradient flow and enables training deep networks.

\item[RMSNorm]
Simplified normalization using only root mean square (no mean centering). Computationally cheaper than LayerNorm.

\item[RoPE (Rotary Position Embedding)]
Position encoding that rotates query and key vectors based on position. Encodes relative positions naturally.

\item[Sampling]
Selecting the next token probabilistically from the model's output distribution. Contrast with greedy decoding.

\item[Scaling Laws]
Empirical relationships between model size, data size, compute, and performance. Chinchilla: optimal compute means N $\propto$ D.

\item[Self-Attention]
Attention where queries, keys, and values all come from the same sequence. The model attends to itself.

\item[Softmax]
Function converting logits to probabilities: $\text{softmax}(x_i) = \frac{\exp(x_i)}{\sum_j \exp(x_j)}$.

\item[Step]
One optimizer update. With gradient accumulation, may involve multiple forward passes.

\item[SwiGLU]
Gated activation function: $\text{SwiGLU}(x) = \text{Swish}(xW_g) \odot (xW_u)$. Outperforms GELU in language models.

\item[Temperature]
Scaling factor for logits before softmax. Higher temperature $\to$ more random; lower $\to$ more deterministic.

\item[Token]
Atomic unit of text for the model. Can be a word, subword, or character depending on tokenization.

\item[Tokenizer]
Algorithm converting text to token IDs and vice versa. \TryLM uses BPE with 16,384 vocabulary size.

\item[Top-k Sampling]
Sampling strategy that only considers the k most probable tokens. Filters out unlikely tokens.

\item[Top-p (Nucleus) Sampling]
Sampling from the smallest set of tokens whose cumulative probability exceeds p. Adapts to probability distribution shape.

\item[Training Loop]
The iterative process: forward pass $\to$ loss computation $\to$ backward pass $\to$ optimizer step.

\item[Transformer]
Neural network architecture based on self-attention. Foundation of modern language models.

\item[UNK (Unknown)]
Special token for out-of-vocabulary inputs. With BPE, rarely needed. In \TryLM: \texttt{<|unk|>}.

\item[Validation]
Evaluating model on held-out data during training. Used to detect overfitting and select best checkpoint.

\item[Vocabulary]
Set of all tokens the model can process. \TryLM uses 16,384 tokens.

\item[Warmup]
Gradually increasing learning rate at the start of training. Helps stability.

\item[Weight Decay]
Regularization adding penalty for large weights: $L_{\text{total}} = L + \lambda \|w\|^2$. AdamW applies this directly to weights.

\item[Weight Tying]
Sharing weights between input embeddings and output projection. Reduces parameters and improves quality.

\end{description}

